\documentclass[11pt,a4paper]{article}
\usepackage{fontspec}
\setmainfont[BoldFont={Noto Sans CJK JP Bold}]{Noto Serif CJK JP}
\setsansfont[BoldFont={Noto Sans CJK JP Bold}]{Noto Sans CJK JP}
\usepackage{amsmath,amssymb,amsthm,mathtools}
\usepackage[margin=24mm]{geometry}
\usepackage{xcolor,booktabs,longtable,array,enumitem,fancyhdr,titlesec}
\usepackage{tocloft}\setlength{\cftsecnumwidth}{3.2em}\setlength{\cftpartnumwidth}{3.2em}
\usepackage[colorlinks=true,linkcolor=blue!50!black,citecolor=blue!50!black,urlcolor=blue!50!black]{hyperref}
\setlength{\parindent}{1em}
\setlength{\parskip}{2pt}
\linespread{1.17}
\emergencystretch=4em
\setlist{itemsep=1pt,topsep=3pt}
\setlength{\tabcolsep}{4pt}

%%% --- Japanese line breaking for XeTeX (no xeCJK available) ---
\makeatletter
\XeTeXinterchartokenstate=1
\newXeTeXintercharclass\cjkc
\newXeTeXintercharclass\cjkclose
\newXeTeXintercharclass\cjkopen
\def\setcjkrange#1#2{\@tempcnta=#1 \@whilenum\@tempcnta<#2 \do{\XeTeXcharclass\@tempcnta=\cjkc \advance\@tempcnta by 1 }}
\setcjkrange{"3000}{"30FF}
\setcjkrange{"4E00}{"9FFF}
\setcjkrange{"FF00}{"FFEF}
\setcjkrange{"2010}{"2027}
\setcjkrange{"25A0}{"25FF}
\setcjkrange{"2190}{"21FF}
\@tempcnta=0
\def\setclose#1{\XeTeXcharclass#1=\cjkclose}
\def\setopen#1{\XeTeXcharclass#1=\cjkopen}
\setclose{"3001}\setclose{"3002}\setclose{"FF0C}\setclose{"FF0E}\setclose{"FF09}\setclose{"300D}\setclose{"300F}\setclose{"3011}\setclose{"FF1A}\setclose{"FF1B}\setclose{"FF01}\setclose{"FF1F}\setclose{"30FC}
\setopen{"FF08}\setopen{"300C}\setopen{"300E}\setopen{"3010}
\def\cjkglue{\hskip0pt plus 0.25pt minus 0.1pt\relax}
\XeTeXinterchartoks\cjkc\cjkc={\cjkglue}
\XeTeXinterchartoks\cjkclose\cjkc={\cjkglue}
\XeTeXinterchartoks\cjkclose\cjkopen={\cjkglue}
\XeTeXinterchartoks\cjkc\cjkopen={\cjkglue}
\XeTeXinterchartoks\cjkclose\cjkclose={}
\XeTeXinterchartoks\cjkc\cjkclose={}
\XeTeXinterchartoks 0 \cjkc={\cjkglue}
\XeTeXinterchartoks \cjkc 0={\cjkglue}
\XeTeXinterchartoks 0 \cjkopen={\cjkglue}
\XeTeXinterchartoks \cjkclose 0={\cjkglue}
\makeatother

%%% --- status tags ---
\newcommand{\tP}{\textcolor{green!45!black}{\textbf{【証明済】}}}
\newcommand{\tC}{\textcolor{blue!55!black}{\textbf{【文献】}}}
\newcommand{\tV}{\textcolor{violet!80!black}{\textbf{【計算で確認】}}}
\newcommand{\tK}{\textcolor{red!70!black}{\textbf{【予想】}}}
\newcommand{\tQ}{\textcolor{orange!85!black}{\textbf{【要確認】}}}

%%% --- theorem-like environments ---
\theoremstyle{definition}
\newtheorem{defn}{定義}
\newtheorem{thm}[defn]{定理}
\newtheorem{prop}[defn]{命題}
\newtheorem{lem}[defn]{補題}
\newtheorem{cor}[defn]{系}
\newtheorem{ex}[defn]{例}
\newtheorem{rem}[defn]{注意}
\newtheorem{obs}[defn]{観察}
\newtheorem{conj}[defn]{予想}
\newcommand{\src}[1]{\par\noindent{\small\textcolor{blue!55!black}{\textbf{[出典]}~#1}}\par}
\renewcommand{\proofname}{証明}

\newenvironment{hyptab}{\par\medskip\noindent{\small\textbf{仮定と、その役割}}\par\nopagebreak\small\noindent\begin{tabular}{>{\raggedright\arraybackslash}p{3.3cm}>{\raggedright\arraybackslash}p{5.5cm}>{\raggedright\arraybackslash}p{5.9cm}}\toprule 仮定 & なぜ必要か & 証明のどこで効くか\\\midrule}{\\\bottomrule\end{tabular}\par\medskip}

%%% --- notation ---
\newcommand{\R}{\mathbb R}\newcommand{\C}{\mathbb C}\newcommand{\Z}{\mathbb Z}
\newcommand{\bT}{\mathbb T}
\newcommand{\Lb}{\overline L}
\newcommand{\Jc}{\mathcal J}
\newcommand{\dbar}{\bar\partial}
\newcommand{\ch}{\mathrm{ch}}\newcommand{\Td}{\mathrm{Td}}
\newcommand{\rk}{\mathrm{rk}}
\newcommand{\HH}{\mathrm{HH}}
\newcommand{\Res}{\mathrm{Res}}
\newcommand{\dlog}{\mathrm{dlog}\,}
\newcommand{\sh}{\sharp}

\pagestyle{fancy}\fancyhf{}
\fancyhead[L]{\small 一般化複素幾何の三種の神器：統合論文}
\fancyhead[R]{\small 研究ノート 2026-11-22（統合版）}
\fancyfoot[C]{\small\thepage}
\titleformat{\section}{\large\bfseries\sffamily}{\thesection}{0.8em}{}
\titleformat{\subsection}{\normalsize\bfseries\sffamily}{\thesubsection}{0.8em}{}


\newcommand{\Q}{\mathbb Q}\newcommand{\N}{\mathbb N}
\newcommand{\Hc}{H_c}
\newcommand{\we}{\wedge}
\setcounter{tocdepth}{1}
\newcommand{\oldname}[1]{\textsf{\small（総合版：#1）}}

\renewcommand{\contentsname}{目次}\renewcommand{\abstractname}{概要}
